%% notes-tex/019-simb-multimodal/sections/3-morphology.tex
%% [[experiments.019-simb-multimodal.scripts.calmorph_variance_analysis]]
%% SOURCE: ceiling + shortlist from
%% experiments/019-simb-multimodal/scripts/morphology_noise_ceiling.py ->
%% results/morphology_noise_ceiling.json and results/morphology_feature_ceiling.csv.
%% Scores from results/round_leaderboards.csv.

\section{Morphology\secstatus{todo}}
\label{sec:morphology}

\subsection{The target\secstatus{todo}}

The \term{morphology strand} predicts the CalMorph phenotype vector of a single-deletion
strain. CalMorph's nominal vocabulary is 501 parameters, which resolve into 281 base
parameters and 220 coefficient-of-variation statistics; the served target is the 281 base
parameters, of which the model scores 278 after three near-degenerate features are dropped.
Source is Ohya 2005 (doi:10.1073/pnas.0509436102), 4,718 deletion mutants plus 122
independent \gene{his3} wild-type replicates.

Raw CalMorph values span about eight orders of magnitude, from $6.7\times10^{-5}$ to
$1.49\times10^{4}$ in absolute mean, so an unnormalized mean-squared-error loss is
dominated by whole-cell size and actin-region size and read $\approx 4.7$ million. The
training path now applies a per-feature $z$-score fit on the training split only, which
brings the loss to order 1 and makes ``predict each feature's mean'' score zero rather than
win. Ohya's own normalization is a per-parameter Box-Cox fit on wild-type data followed by
standardization, and it discards 247 of 501 parameters on a Shapiro-Wilk normality test at
$p \ge 0.5$. That verdict is a normality test over base and CV parameters together, not a
variance floor on the 278 modeled here, so it was not imported.

\subsection{The ceiling was never measured until now, and it is high\secstatus{todo}}

Ohya measures each of the 4,718 deletion strains once and the \gene{his3} wild type 122
independent times, which is exactly what makes reliability estimable without replicating the
mutants. Each strain's value for feature $k$ is one noisy draw $T_k = S_k + E_k$: the
strain's true value plus one sample's measurement noise. The wild-type replicates hold
genotype fixed, so their spread is $E_k$ alone; the mutant panel holds nothing fixed, so its
spread carries $S_k$ and $E_k$ together.

\notesec{Reliability and ceiling are the same measurement on two scales, which is why both
names exist} \term{Reliability} is the fraction of observed across-strain \emph{variance}
that is real signal,
$$\rho_k = \frac{\operatorname{Var}(S_k)}{\operatorname{Var}(T_k)}
  = 1 - \frac{\sigma^2_{\mathrm{WT},k}}{\sigma^2_{\mathrm{mutants},k}},$$
estimated with the sample variance over the 122 replicates and over the 4,718 mutants. The
\term{ceiling} is what a \emph{correlation} can reach. A perfect predictor outputs $S_k$ and
is still scored against the noisy $T_k$, so it correlates at
$$\operatorname{corr}(S_k,\, S_k + E_k)
  = \frac{\operatorname{Var}(S_k)}{\sqrt{\operatorname{Var}(S_k)}\sqrt{\operatorname{Var}(T_k)}}
  = \sqrt{\rho_k} = c_k.$$
The square root is the entire difference: correlation is a ratio of standard deviations
where reliability is a ratio of variances. Since $\rho_k \le 1$ the ceiling always exceeds
the reliability, so the two must never be quoted interchangeably. Over the 278 modeled
features the means are $\rho = 0.444$ and $c = 0.611$, and those are one measurement stated
twice. Compare the ceiling to a reported Pearson; compare the reliability to an $R^2$.

The \term{robust CV} used alongside them is a different axis entirely, measuring how far
deletions move a feature rather than how well it is measured:
$\mathrm{rCV}_k = (Q_{75} - Q_{25}) / |\mathrm{med}|$ over the 4,718 mutants, using
quartiles rather than a standard deviation because the CalMorph features are bounded ratios
and counts with heavy tails.

\begin{table}[htbp]
\centering\small
%% SOURCE: experiments/019-simb-multimodal/results/morphology_noise_ceiling.json
\begin{tabular}{lr}
\toprule
quantity & value \\
\midrule
mutants & 4{,}718 \\
wild-type replicates & 122 \\
modeled features & 278 \\
mean reliability & 0.444 \\
\textbf{mean ceiling over modeled features} & \textbf{0.611} \\
median ceiling & 0.647 \\
features with ceiling $> 0.5$ & 201 of 278 ($72\,\%$) \\
\midrule
best observed (\texttt{vsceij2v}, peak epoch 27 of 65) & \textbf{0.0824} \\
fraction of ceiling realized & $\mathbf{13.5\,\%}$ \\
\bottomrule
\end{tabular}
\caption[]{Morphology has the second-highest ceiling of any strand here and the lowest
fraction of it realized. Nearly three quarters of the modeled features are individually
measurable to $r > 0.5$.}
\label{tab:morph-ceiling}
\end{table}

\subsection{The strand was never really run, on two counts\secstatus{todo}}

Twenty-three runs exist in \file{torchcell_019_morph_v5}, spanning 11 to 121 epochs,
one of them a collapsed constant predictor. The best peaks at epoch 27 of 65. The joint
expression-plus-morphology project is thinner still: 32 runs, none past 87 epochs, best
\file{val/mean/pearson_per_feature} of $0.0505$.

Set against \cref{sec:expression}, where the same architecture on a related target was
still improving at epoch 9,697, a 65-epoch morphology run carries no information about
whether morphology is learnable. The morphology result is not a null. It is an absence of
measurement, and the two must not be recorded the same way.

\notesec{And the training set was a quarter of what exists, in every project} This was
checked against all eight projects that carry a morphology head rather than against one,
because the first version of this claim rested on \file{torchcell_019_morph_v5} alone and
the natural objection is that some other round did train on the full screen. None did.

\begin{table}[htbp]
\centering\small
%% SOURCE: results/round_leaderboards.csv, n_train_supervised per project
\begin{tabular}{lrrr}
\toprule
group & projects & runs & $n_{\mathrm{train}}$ observed \\
\midrule
morphology alone (\file{morph_v2}, \file{_v3}, \file{_v5}) & 3 & 183 & 1{,}161 \\
morphology with expression (\file{expr_morph} .. \file{_v5}) & 5 & 214 & 1{,}161 \\
\midrule
\textbf{total} & \textbf{8} & \textbf{397} & \textbf{1{,}161, no other value} \\
\bottomrule
\end{tabular}
\caption[]{Every morphology run ever launched, across all eight projects that carry a
morphology head, trained on the same 1,161 strains. The Ohya screen holds 4,718 deletions
and 1,440 of them also carry expression; the configs set
\file{require_modalities: [expression_log2_ratio, calmorph]}, which pins every arm to that
intersection, and the 80/20 split takes it to 1,161. For contrast, the same census shows the
expression strand's training set growing across rounds (1,125 to 1,253) and the betaxanthin
rounds differing (3,694 to 4,235), so a constant is a finding here rather than an artifact
of how the census reads.}
\label{tab:morph-ntrain}
\end{table}
\src{experiments/019-simb-multimodal/results/project_census.json}
\src{experiments/019-simb-multimodal/conf/delta_joint_expr_morph_000.yaml}
\src{experiments/019-simb-multimodal/results/fig3_overlap_census.json}

The claim does not rest on the table alone. \file{require_modalities} is set in the config
that every morphology arm inherits, and the overlap census independently gives
\file{ohya_with_expression} $= 1{,}440$ against \file{ohya_unique} $= 4{,}718$. The W\&B
column is corroboration of a restriction that is visible in the configuration.

That restriction is correct \emph{for the auxiliary-task question}, since comparing a joint
model against a morphology-only model on different instance sets would confound the
auxiliary head with a data-quantity difference. It is wrong for the question ``how well can
morphology be predicted'', and the same configs were used for both. \textbf{Morphology has
never been trained on its own full dataset.}

\subsection{Which second label should morphology be paired with\secstatus{todo}}

The overlap census answers this directly, and the answer is not expression.

\begin{table}[htbp]
\centering\small
%% SOURCE: experiments/019-simb-multimodal/results/fig3_overlap_census.json, db_overlap_census
\begin{tabular}{lr}
\toprule
genotypes carrying & count \\
\midrule
morphology (Ohya, total) & 4{,}718 \\
morphology and fitness & \textbf{4{,}220} \\
morphology and expression & 1{,}440 \\
morphology, fitness and expression & 1{,}326 \\
\bottomrule
\end{tabular}
\caption[]{Fitness shares nearly three times as many strains with morphology as expression
does, and covers $89\,\%$ of the morphology screen against expression's $31\,\%$.}
\label{tab:morph-overlap}
\end{table}

Pairing morphology with fitness keeps 4,220 of the 4,718 strains, so the joint arm and the
morphology-only arm can be run on a large common instance set without the restriction above
costing three quarters of the data. Pairing it with expression caps both arms at 1,440.

\notesec{The cap is a consequence of \file{require_modalities}, not a fact about the data}
Both caps come from dropping any genotype that is missing either phenotype. Masking the loss
over absent outputs instead lets a strain with morphology and no expression still contribute
morphology gradient, and a strain with both contribute both. That removes the cap on the
\emph{training} side while leaving the controlled comparison available on the shared subset,
and it is the documented way around this. No run here did it; \file{require_modalities} is
set in the config every morphology arm inherits.

\notesec{The question to ask, stated as one sentence} \textbf{Does adding fitness as a second
label improve morphology prediction at all, on any CalMorph feature or on the average across
them?} That is a yes-or-no question with a correct control (the same architecture, the same
strains, morphology head only), and it is worth answering before any more elaborate
multi-task claim. Report it per feature as well as averaged, because a second label that
helps the 201 well-measured features and does nothing for the rest would be invisible in the
mean.

\notesec{This does not retire the expression-helps-morphology question} That question is
what \file{delta_joint_expr_morph_000.yaml} was built to answer, its control is right, and
its three conditions (\texttt{expr}, \texttt{morph}, \file{expr_morph}) are the correct
decomposition. It has simply never been run past 87 epochs. The intended design is the
broader one: train on all of morphology with masked outputs, then ask whether adding
expression supervision on the 1,440 strains that carry it improves anything. That is a long
and awkward run to tune, which is the practical reason expression is being settled first.

\notesec{What the earlier ``0.05 to 0.14 rise'' was} A flatten-Pearson computed over the
$[B, 281]$ matrix before per-feature normalization, which is dominated by between-feature
scale: large features are large in both prediction and target. After $z$-scoring, the same
smoke run reports $\approx 0$ on that metric. The honest metric is the per-feature
correlation across strains, averaged over features, which is what
\cref{tab:morph-ceiling} scores.

\subsection{A scalar target is a reasonable next step, and this is which scalar\secstatus{todo}}

Ranking single features on reliability alone selects ones that are precisely measured and
nearly constant across mutants, which is the wrong end of the trade. The rank key is
$s_k = c_k \cdot \mathrm{rCV}_k$, the ceiling multiplied by the robust coefficient of
variation, so a feature has to be both measurable and actually moved by deletions.

%% GENERATED by experiments/019-simb-multimodal/scripts/morphology_noise_ceiling.py
%% SOURCE: experiments/019-simb-multimodal/results/morphology_noise_ceiling.json, scalar_shortlist
%% Descriptions are verbatim from Ohya 2005 SI Table 2 "parameter statistics"
%%   $DATA_ROOT/torchcell-library/ohyaHighdimensionalLargescalePhenotyping2005/si/si6.pdf
%%   sha256 037a50927bc5b978ab38b207b264da87550feeeeaaf21bc311f65a5ef13e9638
\begin{table}[htbp]
\centering\footnotesize
%% Ragged-right paragraph columns: a narrow cell holding one long underscored token
%% cannot be justified, and justifying it emits an underfull hbox on every row.
\begin{tabular}{l >{\raggedright\arraybackslash}p{44mm} >{\raggedright\arraybackslash}p{72mm} r r r}
\toprule
feature & CalMorph label & SI description & ceiling & robust CV & score \\
\midrule
\file{A113} & \file{actin_n_ratio} & No actin patch found ratio & $0.873$ & $2.369$ & $2.068$ \\
\file{C124_C} & \file{Medium_bud_ratio} & Ratio of medium bud on stage C & $0.698$ & $0.886$ & $0.618$ \\
\file{A110_A1B} & \file{Actin_f_ratio} & Actin f ratio on stage A1B & $0.637$ & $0.939$ & $0.598$ \\
\file{C125_A1B} & \file{Large_bud_ratio} & Ratio of large bud on stage A1B & $0.784$ & $0.630$ & $0.494$ \\
\file{A105} & \file{actin_a_ratio} & Actin a ratio & $0.655$ & $0.617$ & $0.404$ \\
\file{D212} & \file{nuclear_B_ratio_to_nuclear_AA1BC_cells} & Ratio of B (Nuclear) to A, A1, B and C cells & $0.604$ & $0.642$ & $0.387$ \\
\file{D201} & \file{nuclear_B_ratio} & Ratio of B (Nuclear) & $0.565$ & $0.641$ & $0.362$ \\
\file{A109_A1B} & \file{Actin_e_ratio} & Actin e ratio on stage A1B & $0.476$ & $0.756$ & $0.360$ \\
\file{A103_A1B} & \file{Relative_distance_of_actin_patch_center_from_neck_in_mother} & Relative Distance of actin patch center from neck in mother cell on stage A1B & $0.812$ & $0.431$ & $0.350$ \\
\file{A103_C} & \file{Relative_distance_of_actin_patch_center_from_neck_in_mother} & Relative Distance of actin patch center from neck in mother cell on stage C & $0.534$ & $0.639$ & $0.342$ \\
\file{A110} & \file{actin_f_ratio} & Actin f ratio & $0.700$ & $0.468$ & $0.328$ \\
\file{D215} & \file{nuclear_B_ratio_to_nuclear_A1BC_cells} & Ratio of B (Nuclear) to A1, B and C cells & $0.480$ & $0.578$ & $0.277$ \\
\file{A8-2_C} & \file{Total_brightness_of_actin_region_in_bud} & Actin region brightness in bud on stage C & $0.594$ & $0.449$ & $0.267$ \\
\file{A8-2_A1B} & \file{Total_brightness_of_actin_region_in_bud} & Actin region brightness in bud on stage A1B & $0.595$ & $0.436$ & $0.259$ \\
\file{A121_A} & \file{Maximal_distance_between_patches} & Maximam actin patch length on stage A & $0.492$ & $0.517$ & $0.254$ \\
\file{A8-1_A} & \file{Actin_region_brightness} & Actin region brightness in mother cell on stage A & $0.562$ & $0.444$ & $0.250$ \\
\file{A120_A} & \file{Total_length_of_actin_patch_link} & Total length of actin patch link on stage A & $0.404$ & $0.613$ & $0.248$ \\
\file{D208} & \file{nuclear_B_ratio_to_budded_cells} & Ratio of B (Nuclear) to budded cells & $0.424$ & $0.575$ & $0.244$ \\
\file{A108} & \file{actin_d_iso_ratio} & Actin d (iso) ratio & $0.641$ & $0.375$ & $0.241$ \\
\file{C122} & \file{large_bud_ratio} & Ratio of large bud & $0.693$ & $0.342$ & $0.237$ \\
\bottomrule
\end{tabular}
\caption[]{The twenty best scalar morphology targets, ranked by ceiling multiplied by
robust CV. The list is dominated by bounded class ratios, actin localization class,
bud-size class and nuclear-stage class, which is what a deletion actually moves. Cell
size is the opposite trade and cell shape more so: cell size at the unbudded stage
\file{C11-1_A} has a ceiling of $0.928$ but a robust CV of $0.097$, ranking 101st, and
mother-cell roundness \file{C115_A} has a ceiling of $0.972$ with a robust CV of $0.024$,
ranking 245th of the 275 features with a defined rank key. That last id was described as a
whole-cell axis ratio in the previous revision; SI Table 2 calls it roundness of the mother
cell, and it is a shape parameter rather than a size one. Descriptions are quoted
verbatim from Ohya 2005 SI Table 2, source typography included, and the label column is
the same parameter's name in SI Table 1; the two tables word some parameters
differently, and both wordings are the paper's. In a CalMorph id the
leading letter names the stain the parameter is measured from, C the FITC-ConA cell-wall
image, A the rhodamine phalloidin actin image and D the DAPI nuclear image, and the
suffix names the nuclear stage the cells were binned into, \texttt{\_A}, \texttt{\_A1B}
or \texttt{\_C}, with no suffix meaning all stages pooled. The single-letter classes
inside a description, actin $a$ to $f$ and nuclear $A$ to $F$, are defined by the
drawings in SI Fig.~5B rather than by any text, so those rows rest on that figure.}
\label{tab:morph-scalar}
\end{table}


\notesec{On ``predict cell size first'', which the data does not support} The idea is
attractive because cell size is the CalMorph quantity closest to something intuitive, and it
is the one place the panel gives a clean answer. Cell size is measured superbly and moved
almost not at all.

\begin{table}[htbp]
\centering\small
%% SOURCE: experiments/019-simb-multimodal/results/morphology_noise_ceiling.json, size_features
\begin{tabular}{llrrr}
\toprule
feature & CalMorph description & ceiling & robust CV & rank of 275 \\
\midrule
\file{C11-2_A1B} & Area of daughter cell on stage A1B     & 0.880 & 0.125 & 78 \\
\file{C11-1_A}   & Mother cell size on stage A            & 0.928 & 0.097 & 101 \\
\file{C101_A1B}  & Cell size on stage A1B                 & 0.928 & 0.088 & 118 \\
\file{C103_A}    & Long axis length of mother on stage A  & 0.949 & 0.047 & 209 \\
\file{C115_A}    & Roundness of mother cell on stage A    & 0.972 & 0.024 & 245 \\
\midrule
\file{A113}      & (the shortlist leader, for contrast)   & 0.873 & 2.369 & 1 \\
\bottomrule
\end{tabular}
\caption[]{Cell-size features are the opposite trade from the shortlist: near-perfect
ceilings of $0.88$ to $0.97$, and among the least variable features in the panel. Whole-cell
area at stage A has a robust CV of $0.097$ against a population median of $0.125$, so
\textbf{173 of the 278 features have more spread than whole-cell area}, and its rank key is
23 times smaller than \file{A113}'s. \emph{Correction to revision 2}, which called
\file{C115_A} a whole-cell axis ratio: the CalMorph SI defines it as roundness of the mother
cell, a shape parameter rather than a size one. The one size feature worth noting is bud
area (\file{C11-2_A1B}), which carries $28\,\%$ more spread than mother area, plausibly
because a deletion that stalls the cell cycle changes bud size at fixed nuclear stage.}
\label{tab:morph-size}
\end{table}

\notesec{And on ``predict colony size first''} Colony size is fitness, and fitness is already
the strongest strand in the project, trained on the Costanzo and Kuzmin single-mutant fitness
data. Using it as the morphology warm-up would test the pipeline rather than the phenotype.

\textbf{Recommendation: run the scalar warm-up on \file{A113} or \file{C124_C}}, both
reliable and genuinely variable, and treat the run as a convergence-budget calibration for
the 278-feature target rather than as a result. This is queued behind the expression work
(\cref{sec:directions}) rather than proposed for now.

\begin{keybox}
\textbf{Morphology is the largest unexploited target here, and it has never been measured
properly.} Its ceiling, estimated for the first time, is $0.611$ averaged over 278 features,
with 201 of them individually above $0.5$, which is the second-highest ceiling of any strand.
The best score achieved is $0.0824$, at $13.5\,\%$ of it. But that score comes from a
1,161-strain training set drawn from a 4,718-strain screen, in a run that stopped at epoch
65, and \textbf{every one of the 397 morphology runs across eight projects trained on the
same 1,161 strains} because \file{require_modalities} pins them to the intersection with
expression. \textbf{So $0.0824$ is not a null result about morphology; it is the absence of a
measurement}, and the two must not be recorded the same way. For a second label, fitness
shares 4,220 strains with morphology against expression's 1,440.
\end{keybox}

\notesec{Cannot say} Anything about whether morphology is learnable from genotype, because
no morphology run has passed 121 epochs or seen more than a quarter of the data. Anything
about whether expression and morphology share structure, because the joint project's
longest run is 87 epochs.
